\scanpage{8}\mspage{\handwritten{23}}
\makeatletter\tagsleft@true\makeatother
\[
  0 \handwritten{\longrightarrow} \underline{\mathrm Z}/\handwritten{\delta''}\,\underline{\mathrm Z}
  \handwritten{\longrightarrow \mathord{\tikz[x=1em,y=1ex,baseline=0pt,line width=.045em,line cap=round]{\draw (0,.08)--(.88,.08) (0,1.43)--(.88,1.43) (.19,.08)--(.19,1.43) (.44,.08)--(.44,1.43) (.69,.08)--(.69,1.43);}}'}(Y,\underline A)
  \handwritten{\longrightarrow \mathord{\tikz[x=1em,y=1ex,baseline=0pt,line width=.045em,line cap=round]{\draw (0,.08)--(.88,.08) (0,1.43)--(.88,1.43) (.19,.08)--(.19,1.43) (.44,.08)--(.44,1.43) (.69,.08)--(.69,1.43);}}^{\deleted{\unsure{\prime|*}}}}(Y,\underline B)
  \handwritten{\longrightarrow} 0 \quad . \tag*{(4.38)}
\]

Explicitons deux cas particuliers importants:

\deleted{(\unsure{4.6})}
\uline{Cas particulier} \deleted{\unsure{:|1}}\inserted{(4.6):} $Y$ \uline{est une courbe algébrique sur un corps algébriquement clos}. Alors la \deleted{relation (4.13) est automatiqu\unsure{e.|e. La}} plus grande partie de la discussion précédente devient triviale. La relation (4.13) est automatiquement satisfaite, et on a (4.12 bis) :
\[
  \unsure{\mathbb{B}\mathrm r|\mathrm{Br}}(X)\handwritten{\overset{\sim}{\longrightarrow}}H^{1}(Y,\underline B)\quad ; \tag*{(4.39)}
\]
lorsque de plus les $\handwritten{\delta}_{y}$ sont égaux à $1$, i.e. \handwritten{si} $X$ n'a pas de fibre multiple sur $Y$, alors le deuxième membre s'explicite \deleted{comme étant égal à} par
\[
  0\handwritten{\longrightarrow}\underline{\mathrm Z}/\handwritten{\delta}\underline{\mathrm Z}
  \handwritten{\longrightarrow}H^{1}(Y,\underline{\deleted{\unsure{X|B}}\handwritten{A}})
  \handwritten{\longrightarrow}H^{1}(Y,\underline B)\handwritten{\longrightarrow}0\quad , \tag*{(4.40)}
\]
compte tenu que \deleted{dans} l'on a ici nécessairement $\handwritten{\delta}=\handwritten{\delta'}=\handwritten{\delta''}$ .

\uline{Cas particulier} (4.7): $Y$ \uline{est une courbe algébrique propre sur un corps fini} $k$. Ici encore, (4.13) \deleted{\unsure{et (4.16)}}\deleted{\unsure{sont}}\handwritten{est} automatiquement satisfaite, $k$ étant parfait, \deleted{\unsure{et}}\inserted{\handwritten{ainsi que (4.16) puisque $k(y)$ est de dimension $\leq 1$, d'où $H^{3}(k(y),G_{\mathrm m})=0$; de plus}} on aura $\handwritten{\delta}=\handwritten{\delta'}\handwritten{\deleted{\unsure{=\delta''}}}$, et de même pour les invariants hensélisés et strictement hensélisés. Comme on a encore $\mathrm{Br}(Y)=0$, on trouve ici (4.4):
\[
  0\handwritten{\longrightarrow}\mathrm{Br}(X)\handwritten{\longrightarrow \mathord{\tikz[x=1em,y=1ex,baseline=0pt,line width=.045em,line cap=round]{\draw (0,.08)--(.88,.08) (0,1.43)--(.88,1.43) (.19,.08)--(.19,1.43) (.44,.08)--(.44,1.43) (.69,.08)--(.69,1.43);}}}(Y,\underline B)
  \handwritten{\longrightarrow}\underline{\deleted{\unsure{2|Z}}\handwritten{\mathrm Z}}/\deleted{\unsure{1|\delta}}\handwritten{\Delta}\underline{\mathrm Z}
  \handwritten{\longrightarrow}0\quad ,\qquad\text{où}\quad\handwritten{\Delta\mid\delta}\quad . \tag*{(4.41)}
\]
\deleted{\unsure{compte tenu de|compte tenu du}} D'autre part, si les $\handwritten{\delta}_{y}$ sont égaux à $1$, i.e. si le morphisme déduit de $f:X\handwritten{\longrightarrow}Y$ en passant à la clôture algébrique de $k$ n'a pas de fibre multiple, alors on trouve la suite exacte
\[
  0\handwritten{\longrightarrow}\underline{\mathrm Z}/\handwritten{\delta''}\underline{\mathrm Z}
  \handwritten{\longrightarrow \mathord{\tikz[x=1em,y=1ex,baseline=0pt,line width=.045em,line cap=round]{\draw (0,.08)--(.88,.08) (0,1.43)--(.88,1.43) (.19,.08)--(.19,1.43) (.44,.08)--(.44,1.43) (.69,.08)--(.69,1.43);}}}(Y,\underline A)
  \handwritten{\to \mathord{\tikz[x=1em,y=1ex,baseline=0pt,line width=.045em,line cap=round]{\draw (0,.08)--(.88,.08) (0,1.43)--(.88,1.43) (.19,.08)--(.19,1.43) (.44,.08)--(.44,1.43) (.69,.08)--(.69,1.43);}}}(Y,\underline B)\handwritten{\longrightarrow}0\quad , \tag*{(4.42)}
\]
ce qui est un cas particulier de (4.37), compte tenu de\deleted{s} \inserted{\handwritten{la}} relation\deleted{s} suivante\deleted{s}\unsure{:}
\[
  \handwritten{\bar{\delta}}_{y}=\handwritten{\delta}_{y} \tag*{(4.43)}
\]
pour tout point fermé $y$ de $Y$, qui implique ici $\handwritten{\bar{\delta}}_{y}=1$ d'où $\handwritten{\mathord{\tikz[x=1em,y=1ex,baseline=0pt,line width=.045em,line cap=round]{\draw (0,.08)--(.88,.08) (0,1.43)--(.88,1.43) (.19,.08)--(.19,1.43) (.44,.08)--(.44,1.43) (.69,.08)--(.69,1.43);}}'}(Y,\underline A)=\handwritten{\mathord{\tikz[x=1em,y=1ex,baseline=0pt,line width=.045em,line cap=round]{\draw (0,.08)--(.88,.08) (0,1.43)--(.88,1.43) (.19,.08)--(.19,1.43) (.44,.08)--(.44,1.43) (.69,.08)--(.69,1.43);}}}(Y,\underline A)$.

La démonstration de\deleted{s} \handwritten{l'}égalité\deleted{s} (4.43), qui est une conséquence facile du
\reversemarginpar\margin{\handwritten{\unsure{x|$\times$}}}\normalmarginpar
théorème de \deleted{Lang}\handwritten{LANG} \inserted{\handwritten{[24] $H^{1}(k,G)=0$}} pour le corps fini $k(y)$ et le schéma de Picard \deleted{\unsure{con}}\handwritten{con}nexe $\handwritten{G=}\underline{\mathrm{Pic}}^{0}_{X_{y}/y}$, est laissé\handwritten{e} au lecteur comme un exercice plaisant et délectable.
% archive sheet number: 7
\clearpage
